\documentclass[12pt]{article}
\usepackage[utf8]{inputenc}
\usepackage[T1]{fontenc}
\usepackage[margin=1in]{geometry}
\usepackage{setspace}
\usepackage{xcolor}
\usepackage[colorlinks=true, urlcolor=blue!60!black, linkcolor=blue!60!black]{hyperref}
\setlength{\parindent}{0em}
\setlength{\parskip}{0.6em}
\renewcommand{\baselinestretch}{1.0}
\begin{document}
\begin{center}
{\LARGE \textit{Arizona v. United States}, 567 U.S. 387 (2012)\par}
\vspace{0.5em}
{\large Case Brief --- Two-Page Edition\par}
\end{center}

\section{Memory Jogger}
Arizona's S.B.~1070: state crimes for registration failures and unauthorized work, and warrantless arrests for removability, are preempted; the ``show me your papers'' check survives a facial challenge (\href{https://supreme.justia.com/cases/federal/us/567/387/}{Justia}).

\section{Detailed Case Facts}
Arizona enacted S.B.~1070 to pursue ``attrition through enforcement.'' The United States challenged four provisions on preemption grounds only: Section 3 (state crime of failing to carry registration papers), Section 5(C) (state crime for an unauthorized alien to seek or perform work), Section 6 (warrantless arrest on probable cause of a removable offense), and Section 2(B) (status checks with ICE during lawful stops).

\section{Procedural History}
The district court preliminarily enjoined all four provisions; a divided Ninth Circuit affirmed.

\section{Judicial Votes}
5--3 (Kagan recused). Kennedy wrote for Roberts, Ginsburg, Breyer, and Sotomayor. Scalia and Thomas would have upheld all four provisions; Alito agreed on Sections 2(B) and 3 but would have upheld Sections 5(C) and 6.

\section{Holding}
Sections 3, 5(C), and 6 are preempted: Congress occupies the field of alien registration, and the other two obstruct Congress's chosen balance and federal enforcement discretion. Section 2(B) is not facially preempted before state courts construe it, provided it does not prolong detention solely to check status.

\section{Analysis of Opinions}
\textbf{Majority:} Under \textit{Hines v. Davidowitz}, even parallel state registration penalties are barred. IRCA deliberately chose civil, not criminal, consequences for unauthorized employees, and Section 6 gave state police more arrest power than federal officers have. Congress encourages information sharing, so Section 2(B) could stand pending construction. \textbf{Scalia:} States have inherent sovereign power to exclude, and executive nonenforcement cannot preempt laws enforcing federal statutes. \textbf{Thomas:} only enacted text, not implied purposes, preempts. \textbf{Alito:} the majority gave ``short shrift'' to the presumption against preemption.

\section{Examples: Future Applications}
\textbf{Same-Side:} a state crime of unlawful entry with state-ordered removal is preempted. \textbf{Opposite-Side:} a neutral identity-theft prosecution for W-4 fraud is not, as \textit{Kansas v. Garcia} later held. \textbf{Fence-Sitter:} a mandate that jails honor 48-hour ICE detainers may be ``cooperation,'' but holding people past release raises the Fourth Amendment concerns the Court flagged.

\section{Critique}
Civil-rights scholars faulted the Court for treating the case as a pure federal-state contest and ignoring racial profiling under Section 2(B). Textualists argue that implied preemption from congressional silence departs from the Supremacy Clause and turns executive discretion into a preemption shield.

\section{Key Quotations}
\begin{quote}
``The Government of the United States has broad, undoubted power over the subject of immigration and the status of aliens.''
\end{quote}
\noindent --- Justice Kennedy, for the Court, 567 U.S. at 394.

\end{document}
